\documentclass{article}

% NeurIPS 2026 style.
\usepackage{neurips_2026}

\usepackage[utf8]{inputenc}
\usepackage[T1]{fontenc}
\usepackage{hyperref}
\usepackage{url}
\usepackage{booktabs}
\usepackage{amsfonts}
\usepackage{amsmath}
\usepackage{amssymb}
\usepackage{nicefrac}
\usepackage{microtype}
\usepackage{graphicx}
\usepackage[table]{xcolor}
\usepackage{algorithm}
\usepackage{algorithmic}
\usepackage{subcaption}
\usepackage{multirow}
\usepackage{wrapfig}

% The NeurIPS style defines \AND for author blocks; algorithmic also defines
% \AND inside pseudocode. The paper uses only \And in the author field, so we
% undefine the unused all-caps variant before algorithmic environments.
\makeatletter
\let\AND\@undefined
\makeatother

\graphicspath{{figures/}{../outputs/standard/analysis/}{../outputs/vla/}}

\newcommand{\attnres}{\textsc{AttnRes}}
\newcommand{\reskip}{\textsc{ReSkip}}
\newcommand{\reloop}{\textsc{ReLoop}}
\newcommand{\retrofit}{\textsc{AR-Retrofit}}
\newcommand{\routeA}{\textsc{Route-A}}
\newcommand{\routeB}{\textsc{Route-B}}
\newcommand{\routeC}{\textsc{Route-C}}
\newcommand{\E}{\mathbb{E}}
\newcommand{\R}{\mathbb{R}}
\newcommand{\softmax}{\text{softmax}}
\newcommand{\todo}[1]{\textcolor{red}{[TODO: #1]}}

\title{ReSkip: Brain-Inspired Adaptive Depth for Pretrained Transformers\\via Attention Residuals}

\author{
  Anonymous Author(s) \\
  Affiliation withheld \\
  \texttt{anonymous@submission.invalid}
}

\begin{document}

\maketitle

\begin{abstract}
Different inputs deserve different amounts of computation---a System-1 / System-2 dichotomy~\citep{kahneman2011thinking} with concrete neural correlates: shallow feedforward processing for easy stimuli, recurrent and deeper compute for hard ones~\citep{lamme2000distinct, kar2019recurrent, kietzmann2019recurrence, spoerer2020recurrent}. Recent test-time compute scaling for LLMs realises this intuition along the \emph{token} axis (longer reasoning traces for harder problems)~\citep{openai2024o1, deepseek2025r1, snell2024scaling, wei2022chain}, but every token in such a trace still traverses every layer of the transformer; the complementary \emph{depth} axis---which blocks to invoke per input---is largely under-exploited, and existing methods (early exit, mixture-of-depths, layer pruning) bolt on a separate decider rather than producing the depth-allocation signal as a property of the forward pass itself. We use \emph{Attention Residuals} (\attnres{})~\citep{chen2026attnres}---a residual that attends over previous block outputs---as the intrinsic handle: each block's input arrives with a routing weight saying how much each predecessor contributed, and we read this weight as a pre-execution skip signal (\reskip{}). Because \attnres{} is an architectural change, the practical bottleneck is installation; our main contribution, \retrofit{}, is a $\gamma$-gated residual-injection fine-tune that converts any pretrained decoder transformer into an \attnres{}-capable model with the base frozen, well under $1\%$ new parameters, and tens of GPU-minutes per billion. On Qwen3-VL-2B and 4B the retrofit strictly improves the base on most \texttt{lmms-eval} benchmarks; a LIBERO VLA warm-started from the retrofit beats a matched pure-OFT baseline at both scales, and \reskip{} on top of the trained policy beats the no-skip backbone on the 4-suite mean. Under \texttt{torch.compile} the pipeline runs at iso-cost with the compiled base---accuracy gains are obtained at the same fast forward, not against it. To our knowledge this is the cheapest existing path from an off-the-shelf transformer to a depth-axis adaptive-compute model, complementary to token-axis test-time-compute approaches.
\end{abstract}

%----------------------------------------------------------------------
\section{Introduction}
\label{sec:intro}
%----------------------------------------------------------------------

\paragraph{Different inputs deserve different amounts of computation.}
This intuition has two well-developed roots. Cognitively, Kahneman's dual-process framework~\citep{kahneman2011thinking} contrasts \emph{System~1} (fast, automatic, low-effort) with \emph{System~2} (slow, deliberate, effortful), and posits that a competent agent allocates which system handles a problem based on its difficulty. Neurally, this dichotomy has direct correlates in cortical processing: easy stimuli are recognised through shallow feedforward sweeps, while hard ones recruit recurrent and deeper circuits~\citep{lamme2000distinct, kar2019recurrent, dicarlo2012does}. The closest computational-neuroscience analogue is explicit input-dependent depth allocation: recurrent vision models capture human representational dynamics~\citep{kietzmann2019recurrence}, and confidence-thresholded recurrence spends more steps on harder images to trade speed for accuracy~\citep{spoerer2020recurrent}. This also has a natural depth interpretation: finite-time recurrence can be unrolled into feedforward depth, but recurrence reuses the same substrate to obtain flexible effective depth~\citep{vanbergen2020going}. At the mechanism level, the routing decision is plausibly intrinsic to cortical circuitry~\citep{bastos2012canonical, buschman2007topdown}; dendritic integration provides a cellular route by which feedforward and contextual feedback inputs can be associated inside the neuron itself~\citep{larkum2013cellular}, not by an external supervisor. Both perspectives describe the same competence: \emph{the agent decides its own depth, per input}.

\paragraph{Modern AI does System~2, but only on the token axis.}
Recent work on large language models operationalises System-2-style allocation, but along a single axis: \emph{tokens}. Test-time compute scaling~\citep{openai2024o1, deepseek2025r1, snell2024scaling}, chain-of-thought prompting~\citep{wei2022chain}, and self-thought methods~\citep{zelikman2024quietstar} all expose more compute on hard problems by emitting more reasoning tokens before answering. This is one valid form of System~2: longer chains of identical-depth processing. A second, complementary axis is largely under-exploited: \emph{depth}. Every token in a long reasoning trace still goes through every layer of the underlying transformer, regardless of whether that token is a routine connector or a load-bearing inferential step. The cognitive analogy suggests both axes should be available; current architectures expose only one.

\paragraph{Existing depth-axis methods bolt on the decider.}
Early-exit classifiers~\citep{schuster2022confident, elbayad2020depth} add auxiliary heads per layer; Mixture-of-Depths~\citep{raposo2024mixture} learns a token router with a capacity-balancing loss; layer-pruning~\citep{gromov2024unreasonable, men2024shortgpt} is post-hoc and static; Universal Transformers~\citep{dehghani2018universal} share weights with ACT-style halting~\citep{graves2016adaptive}. In each case a separate component decides whether the rest of the network should run---the dual-process gating is \emph{added on top}, not produced \emph{by} the forward pass. Cognitive and neural System-2 control, in contrast, is not a separate module; it is a property of the substrate.

\paragraph{Attention Residuals provide exactly such an intrinsic handle.}
\attnres{}~\citep{chen2026attnres} replaces the fixed scalar-$1$ residual with a learned softmax-attention over previous block outputs; the routing weights $\alpha_{i\to n}$ are input-dependent, competitively normalised, and---crucially---computed \emph{before} block $n$ runs, as part of forming its input. They emit a per-block depth-allocation signal as a side-effect of the network's own forward, not as a separate head. A block whose downstream consumers rarely attend to it is, by construction, a block whose computation is locally redundant for that input.

\paragraph{The bottleneck: AttnRes is an architectural change, and changing the architecture is expensive.}
Every existing \attnres{} model is pretrained from scratch with the modified residual; reproducing this at the 2--7B scale that real applications need costs tens of thousands of GPU-hours (Appendix~\ref{app:motivation}). The community has already invested in strong standard-residual transformers---Qwen3-VL, LLaVA-OneVision, InternVL---that no group will retrain just to acquire a depth-axis routing signal. If depth-axis System~2 is to be available for the standard pretrained models the community already trusts, we need a way to install it through a short fine-tune.

\paragraph{Contributions.}
\begin{itemize}
\item \textbf{\reskip{}} (\S\ref{sec:reskip}). The \attnres{} routing weights drive a per-input layer-skip rule with no auxiliary head; the importance--ablation disconnect requires combining both signals to pick eligible positions.
\item \textbf{\retrofit{}} (\S\ref{sec:retrofit}, main contribution). A $\gamma$-gated residual injection installs \attnres{} into a frozen pretrained transformer through a single short fine-tune ($\sim$$0.7\%$ new parameters, identity-at-init, tens of GPU-minutes per billion), with every block converging to $\gamma_n{=}1$.
\item \textbf{VLA application} (\S\ref{sec:vla}). A LIBERO policy warm-started from the retrofit improves the matched pure-OFT baseline under our deployment-style reporting convention and remains positive under mean-of-seeds sensitivity; \reskip{} on top of the trained policy improves the no-skip policy in the same calibration sweep.
\item \textbf{Iso-cost} (\S\ref{sec:exp_retrofit}). Under \texttt{torch.compile} the retrofit forward measures $1.029\times$ \texttt{base\_compiled}: $+$accuracy at the same fast forward, not a speed/accuracy trade-off.
\end{itemize}


%----------------------------------------------------------------------
\section{Method}
\label{sec:method}
%----------------------------------------------------------------------

\subsection{Background: Attention Residuals}
\label{sec:background}

A standard transformer uses fixed scalar-$1$ combination, $\mathbf{x}_l = \mathbf{x}_{l-1} + f_l(\mathbf{x}_{l-1})$, emitting no routing signal. \attnres{}~\citep{chen2026attnres} replaces this with a learned attention over depth. Each block $l$ maintains a pseudo-query $\mathbf{w}_l \in \R^d$; keys are projections $\mathbf{k}_i = W_K \mathbf{h}_i$ of earlier block outputs $\mathbf{h}_i$, and
\begin{equation}
  \alpha_{i \to l} = \frac{\exp(\mathbf{w}_l^\top \mathbf{k}_i / \sqrt{d})}{\sum_{j=0}^{l-1} \exp(\mathbf{w}_l^\top \mathbf{k}_j / \sqrt{d})},
  \qquad
  \mathbf{x}_l = \sum_{i=0}^{l-1} \alpha_{i \to l}\,\mathbf{h}_i .
  \label{eq:attnres}
\end{equation}
\textbf{Two-phase execution.} Computing the input to block $l$ requires only $\{\mathbf{h}_i\}_{i<l}$ and $\mathbf{w}_l$---both available \emph{before} $f_l$ runs. We call the pre-execution attention computation \emph{phase~1} and $f_l$ itself \emph{phase~2}. Phase~1 is the routing signal we exploit. Block-\attnres{} groups layers into $N$ blocks and applies the softmax at the block level (the granularity used throughout). Algorithm~\ref{alg:two_phase} (Appendix~\ref{app:algorithm}) gives the full two-phase forward with \reskip{}'s skip check folded in.

\subsection{\reskip{}: AttnRes-guided layer skipping}
\label{sec:reskip}

\reskip{} turns the phase-1 signal into a per-input skip rule. Each block~$n$ has two offline scores: an \emph{\attnres{} importance} $I(n) = \max_{l > n}\,\E[\alpha_{n \to l}]$ (how often block $n$ is referenced; forward passes only) and an \emph{ablation impact} $A(n) = \mathrm{PPL}(\text{$n$ removed}) / \mathrm{PPL}(\text{full})$ (is block $n$ irreplaceable; $N{-}2$ ablated forwards). At inference, for each $n$ we read the mean weight $w_{\text{recent}}(n)$ that position $n$'s router places on the immediate predecessor and apply
\begin{equation}
  \text{skip block $n$} \iff
  w_{\text{recent}}(n) > \tau_n
  \;\text{and}\;
  n \in \mathcal{P}
  \;\text{and}\;
  \textstyle\sum_{m<n} \mathbb{1}[\text{skipped}_m] < M_{\max} ,
  \label{eq:skip_rule}
\end{equation}
with $\tau_n$ calibrated per position as the $q$-th empirical quantile of $w_{\text{recent}}(n)$ on $32$ held-out long-context sequences ($q{=}0.85$ default). The skip decision is a scalar comparison on a quantity already computed in phase~1; no auxiliary network, no train-time changes, no architectural modification (cf.\ CALM / MoD / pruning / LayerSkip in Appendix~\ref{app:method_comparison}).

\paragraph{Importance--ablation disconnect.}
$I(n)$ and $A(n)$ disagree on a case-by-case basis: block~$2$ on our 340M from-scratch \attnres{} (Figure~\ref{fig:reskip_routing}) has middle-of-the-pack importance ($I{=}0.52$) yet is \emph{catastrophic} to remove ($A{=}8.1\times$ PPL); block~$4$ has the lowest ablation impact yet higher $I$ than block~$5$. $\mathcal{P}$ should therefore use both signals---low $I$ so skip \emph{triggers often}, low $A$ so it is \emph{safe}. Picking by either alone is strictly worse: $\{5\}$ alone (low-$I$) gives $1.14\times$; $\{3,5\}$ combined gives $1.19\times$ at PPL slightly below full depth. Full position sweep in Appendix~\ref{app:reskip_position_ablation}.

\begin{figure}[t]
\centering
\includegraphics[width=0.95\textwidth]{reskip_routing_importance_vs_ablation.pdf}
\caption{\textbf{\attnres{} routing weights do not predict ablation impact (340M).} (a)~Learned block-level $\alpha_{n \to l}$. (b)~$I(n)$ (blue, left axis) and $A(n)$ (orange, right axis). \reskip{} needs both signals.}
\label{fig:reskip_routing}
\end{figure}

\paragraph{Validation at 340M from-scratch.}
On a 340M block-\attnres{} transformer ($d{=}1024$, $L{=}24$, $N{=}8$ blocks; FineWeb-Edu 100BT~\citep{penedo2024fineweb}; \texttt{lm-eval-harness}~\citep{eval-harness}), \reskip{} at $\mathcal{P}{=}\{3,5\}, M_{\max}{=}2, q{=}0.85$ reaches \textbf{$1.19\times$ wall-clock at zero benchmark drop} on LAMBADA~\citep{paperno2016lambada}, HellaSwag~\citep{zellers2019hellaswag}, and ARC-E/C~\citep{clark2018arc} (Table~\ref{tab:reskip_benchmark}; Pareto and latency curves in Appendix~\ref{app:reskip_results}). Skip count varies $0$--$2$ per forward (mean $0.88$ over $48$ batches), confirming genuine input-dependence rather than a fixed-block removal. A 110M cross-scale replication on the same FineWeb-Edu recipe reproduces the zero-degradation pattern at $q{=}0.97, M_{\max}{=}1$ (Appendix~\ref{app:phase1}). \emph{This validates that the mechanism is real when \attnres{} is trained in. The harder question---how to acquire it in an already-trained model---is \S\ref{sec:retrofit}.}

\begin{table}[t]
\caption{\textbf{\reskip{} on 340M from-scratch \attnres{}}: dynamic skip preserves every metric while running $1.19\times$ faster end-to-end at sequence length $8192$.}
\label{tab:reskip_benchmark}
\centering\small
\begin{tabular}{lccccc}
\toprule
Configuration & LAMBADA acc & LAMBADA ppl & HellaSwag & ARC-E & ARC-C \\
\midrule
Full-depth & 0.4056 & 20.20 & 0.4607 & 0.5438 & 0.3012 \\
\reskip{} ($\{5\}, M{=}1, q{=}0.95$) & \textbf{0.4056} & \textbf{20.20} & \textbf{0.4607} & \textbf{0.5438} & \textbf{0.3012} \\
\reskip{} ($\{3,5\}, M{=}2, q{=}0.85$) & \textbf{0.4056} & \textbf{20.20} & \textbf{0.4607} & \textbf{0.5438} & \textbf{0.3012} \\
\bottomrule
\end{tabular}
\end{table}

\subsection{\retrofit{}: installing AttnRes into a pretrained transformer}
\label{sec:retrofit}

\paragraph{Goal.} Take a transformer pretrained in the standard-residual regime and produce an \attnres{}-capable model that (P1) preserves benchmark quality, (P2) emits a routing $\alpha$ informative enough to drive \reskip{}, and (P3) trains in a single short fine-tune over off-the-shelf SFT data, with the base frozen.

\paragraph{$\gamma$-gated residual injection.}
We organise the $L$ decoder layers into $N$ \emph{blocks} of $L/N$ layers each (Qwen3-VL-2B: $L{=}28, N{=}14$). For every block $n\ge 1$ we add an \attnres{} router $\mathbf{w}_n$, a tiny adapter $A_n$ (down-up bottleneck of rank $r$, SiLU), and a scalar gate $\gamma_n$:
\begin{equation}
  r_n = \sum_{i=0}^{n-1} \alpha_{i \to n}\,h_i, \qquad
  x_n = h_{n-1} + \gamma_n\,A_n(r_n - h_{n-1}), \qquad
  h_n = \mathrm{Block}_n(x_n).
  \label{eq:gamma_gate}
\end{equation}
$\mathrm{Block}_n$ is the pretrained block, verbatim; \attnres{} enters \emph{between} blocks, as a correction. The retrofit parameters $\{\mathbf{w}_n, \gamma_n, A_n\}$ amount to $\sim$$15$M ($\sim$$0.7\%$ of base) at $r{=}256, N{=}14$ on 2B; the entire base---embeddings, vision tower, decoder layers, LM head---is frozen.

\paragraph{Identity at init.}
With $\gamma_n{=}0$ we have $x_n = h_{n-1}$ exactly, so the forward at step $0$ is bit-identical to the pretrained model. The adapter up-projection is small-random ($\mathcal{N}(0,0.02^2)$) so that $\partial\mathcal{L}/\partial\gamma_n \ne 0$ at init, avoiding gradient deadlock. We then ramp $\gamma_n$ via a linear curriculum $0\to 1$ over the first $30$--$50\%$ of training; every block converges to $\gamma_n{=}1$, leaving the retrofit \emph{structurally pure \attnres{}} (every block consumes the routed sum, not the standard residual). Three cleaner-looking alternatives we tried first---an observer-only head, an interpolation $\beta\to 1$ at the block input, and an informed-init pseudo-query with temperature anneal---all fail (Appendix~\ref{app:route_ablations}).

\paragraph{Skip integrates with the same path.}
\reskip{}'s phase-1 decision reads the already-computed $\alpha$ and, when the rule triggers, short-circuits the block: $h_n\!\leftarrow\!x_n$. To make this consistent with \texttt{use\_cache=True}, on a skipped block we still run the per-layer K/V-producing slice ($\texttt{LayerNorm}\!\to\!\texttt{k/v\_proj}\!\to\!\texttt{RoPE}\!\to\!\texttt{cache.update}$), which costs ${\approx}5$--$10\%$ of a full-layer forward, and skip \texttt{q\_proj} / attention / \texttt{o\_proj} / MLP. The combined \texttt{skip\,+\,use\_cache=True} path matches the cacheless path in last-position argmax (Appendix~\ref{app:kv_equiv}).

\paragraph{Training objective.}
We minimise
\begin{equation}
  \mathcal{L} = \mathcal{L}_{\text{CE}}(\text{full path}) + \lambda_{\text{kl}}\,\mathcal{L}_{\text{KL}}(\text{skip path}\,\|\,\text{teacher}) + \lambda_{\text{ent}}\,\mathcal{L}_{\text{ent}}(\alpha) ,
  \label{eq:retrofit_loss}
\end{equation}
where the full-path CE is computed on assistant tokens, while the skip-branch KL samples one block at random per step, runs the forward with it skipped, and pulls the resulting logits toward a frozen pretrained teacher (so the surrogate $x_n$ is useful both when the block runs and when it is skipped). $\mathcal{L}_{\text{ent}}$ is implemented with a negative sign in the loss, i.e., a light entropy-maximising prior against premature collapse of $\alpha$. Default $\lambda_{\text{kl}}{=}1$, entropy weight $0.02$, lr $10^{-3}$ on retrofit params, cosine schedule with $100$-step warmup. Hyperparameters and the v1 / v3 data mixes used throughout are listed in Appendix~\ref{app:retrofit_hparams}.

\paragraph{Why this objective is not ordinary SFT.}
The CE term keeps the full path useful for the downstream task, but by itself it does not train the skipped surrogate to replace a real block. The skip-branch KL supplies that constraint by tying one skipped-block forward per step to the frozen teacher. Conversely, teacher imitation alone would merely reproduce the base model. The weak entropy prior prevents early one-source collapse while still allowing the router to sharpen later. This three-term structure is what makes the retrofit identity-preserving at step $0$, task-improving after training, and skip-ready at inference.

%----------------------------------------------------------------------
\section{Experiments: \retrofit{} on Qwen3-VL-2B and 4B}
\label{sec:experiments}
\label{sec:exp_retrofit}
%----------------------------------------------------------------------

\paragraph{Setup.}
We retrofit \textbf{Qwen3-VL-2B}~\citep{bai2025qwen3vl} ($L{=}28, d{=}2048$) at $N{=}7$ blocks of $L{=}4$ and \textbf{Qwen3-VL-4B} ($L{=}36, d{=}2560$) at $N{=}9$ blocks of $L{=}4$ (block-partition sweep in Appendix~\ref{app:block_partition}). $L{=}4$ is the canonical partition for both scales: it sits on the accuracy plateau alongside $L{=}2$ and halves the per-token router count, which is the dominant factor in the inference-cost result below. All base parameters are frozen; $\sim$$15$M ($2$B) / $\sim$$23.7$M ($4$B) retrofit parameters at $r{=}256$ train via Eq.~\ref{eq:retrofit_loss}. We compare two data mixes: \textbf{v1} ($50$/$50$ UltraChat-$200$k~\citep{ding2023ultrachat} + LLaVA-Instruct-VSFT~\citep{liu2023llava}, $5$k steps; $\sim$$22$ GPU-min on $2$B), used as a LoRA-comparable narrow-mix control; and \textbf{v3} ($60\%$ LLaVA-OneVision-Data~\citep{li2024llavaonevision} + $20\%$ UltraChat + $10\%$ NuminaMath + $10\%$ OpenThoughts, $10$k steps), the canonical mix used for every result we report. Evaluation: LAMBADA~\citep{paperno2016lambada}/HellaSwag~\citep{zellers2019hellaswag} for text; six \texttt{lmms-eval} benchmarks for VLM (MMBench~\citep{liu2023mmbench}, MMMU~\citep{yue2024mmmu}, MMStar~\citep{chen2024mmstar}, AI2D~\citep{kembhavi2016ai2d}, OCRBench~\citep{liu2024ocrbench}, RealWorldQA~\citep{grok2024realworldqa}); text-side \reskip{} eligible sets selected from per-block static removal, and VLA eligible sets selected from action-drift sweeps on the trained policy (Appendices~\ref{app:block_removal},\,\ref{app:vla_pareto}). Identity-at-init verified ($\gamma{=}0$ reproduces base; max $|\Delta\text{logits}|=0.375$ bf16, $100\%$ argmax agreement; Appendix~\ref{app:identity}).

\begin{table}[t]
\caption{\textbf{\retrofit{} on Qwen3-VL-2B/4B at the canonical $L{=}4$ v3 recipe} (\texttt{lmms-eval} full splits for VLM, $n{=}2000$ for LAMBADA/HellaSwag). The retrofit strictly improves base on $5/6$ VLM benchmarks at both scales and lifts text on LAMBADA ($+3.3$pp on $2$B; $+8.7$pp on $4$B). MMStar overall holds; the math sub-category (Appendix~\ref{app:mmstar_subcat}) jumps $+7.9$pp on $2$B / $+3.9$pp on $4$B---the largest single cell across the suite, consistent with a router that lifts deliberate-reasoning paths (\S\ref{sec:exp_retrofit}). Parameter-matched LoRA baselines fail to replicate the text gain ($-0.6$pp LAMBADA over four runs; Appendix~\ref{app:lora_baselines}). All accuracy results in this table are obtained on the iso-cost forward of \S\ref{sec:exp_retrofit} ``Inference cost''.}
\label{tab:retrofit_main}
\centering\small
\setlength{\tabcolsep}{4pt}
\begin{tabular}{lcccccccc}
\toprule
& \multicolumn{2}{c}{Text} & \multicolumn{6}{c}{VLM (\texttt{lmms-eval} full split)} \\
\cmidrule(lr){2-3}\cmidrule(lr){4-9}
Configuration & LAMBADA & HellaSwag & MMBench & MMMU & MMStar & AI2D & OCRBench & RWQA \\
\midrule
Qwen3-VL-2B (base)                    & 0.532 & 0.506 & 75.77 & 0.414 & 0.536 & 0.736 & 0.772 & 0.648 \\
$+$ \retrofit{} v3 ($L{=}4$, $10$k)   & \textbf{0.5650} & 0.500   & \textbf{78.87} & \textbf{0.432} & 0.536 & \textbf{0.758} & \textbf{0.814} & \textbf{0.661} \\
~~~~$\Delta$ vs.\ base                & \textbf{+3.3}   & $-0.6$  & \textbf{+3.10} & \textbf{+1.8} & $0.0$  & \textbf{+2.2}  & \textbf{+4.2}  & \textbf{+1.3}  \\
\midrule
Qwen3-VL-4B (base)                    & 0.576 & 0.562 & 83.33 & 0.490 & 0.624 & 0.819 & 0.819 & 0.715 \\
$+$ \retrofit{} v3 ($L{=}4$, $10$k)   & \textbf{0.6625} & 0.5515 & \textbf{85.22} & \textbf{0.521} & \textbf{0.632} & \textbf{0.825} & \textbf{0.824} & \textbf{0.718} \\
~~~~$\Delta$ vs.\ base                & \textbf{+8.7}   & $-1.1$ & \textbf{+1.9}  & \textbf{+3.1}  & \textbf{+0.8}  & \textbf{+0.6}  & \textbf{+0.5}  & \textbf{+0.3}  \\
\bottomrule
\end{tabular}
\end{table}


\paragraph{Quality.}
The canonical retrofit tells a consistent story across both scales (Tab.~\ref{tab:retrofit_main}): strictly above base on five of six VLM benchmarks at $2$B and at $4$B, with the largest gain on the MMStar math subcategory ($+7.9$pp at $2$B, $+3.9$pp at $4$B; full $6$-cell decomposition in Appendix~\ref{app:mmstar_subcat}). The win concentrates on \emph{deliberate reasoning over images} rather than perception, consistent with a router that lifts the slow-thinking path. Text quality also lifts: $+3.3$pp LAMBADA / $-16\%$ ppl on $2$B and $+8.7$pp / $-32\%$ on $4$B; HellaSwag is within $\pm 1.1$pp at both scales. Every block converges to $\gamma_n{=}1$, so the retrofit is structurally pure \attnres{}. Parameter-matched LoRA baselines on the same data produce $-0.6$pp LAMBADA over four runs ($5.0$pp behind the retrofit, Appendix~\ref{app:lora_baselines}): the gain comes from \emph{the \attnres{} routing structure}, not just from extra trainable parameters.

\paragraph{Attribution controls.}
The result is not simply ``small adapter fine-tuning helps Qwen3-VL''. LoRA at the same parameter scale and twice the narrow-mix step count averages $0.526$ LAMBADA versus $0.532$ for the frozen base and $0.576$ for the matching v1 \retrofit{} run. Observer-only routing learns an $\alpha$ signal that never feeds back into the forward, and $\gamma$-free / interpolation variants destabilise the frozen backbone before the router stabilises. The converged recipe therefore depends on the identity-preserving bridge, the real forward-path \attnres{} signal, and the skip-aware surrogate loss together.

\paragraph{\reskip{} on the retrofit.}
With per-scale eligible sets ($2$B: $\mathcal{P}{=}\{1,4\}$; $4$B: $\mathcal{P}{=}\{1,2\}$) and per-block thresholds calibrated as the $q$-th quantile of $w_{\text{recent}}$ on a $32$-prefix held-out set, dyn-skip preserves LAMBADA at the conservative end ($q{=}0.85$ holds within $1.0$pp of no-skip; full $(q,M)$ sweep in Appendix~\ref{app:retrofit_pareto}). The same Pareto shape reappears on action prediction (\S\ref{sec:vla}): there \reskip{} at $q{=}0.99$ \emph{beats} the no-skip backbone on the LIBERO 4-suite mean, validating that the depth-axis allocation transfers from token-level to action-level decoding under the same threshold-calibration protocol.

\paragraph{Inference cost: iso-cost under \texttt{torch.compile}.}
Two architectural choices put the accuracy results on the same fast forward path. (i) The canonical $L{=}4$ partition halves the per-token router count vs $L{=}2$, closing $\sim$$70\%$ of the eager-vs-eager gap on its own ($1.39\times \to 1.12\times$, seq~$2048$, cache, H100/bf16; Tab.~\ref{tab:retrofit_latency}). (ii) Both base and retrofit wrapped in \texttt{torch.compile} (\texttt{max-autotune}, \texttt{dynamic=False}) fuse the small router gemms; the retrofit benefits more because it has more router-shaped kernels. The retrofit measures $\mathbf{1.029\times}$ \texttt{base\_compiled} and $\mathbf{0.864\times}$ of the stock-HF \texttt{base\_eager} forward---i.e. \emph{\retrofit{} $+$ \texttt{compile} runs $14\%$ faster than what most Qwen3-VL-$2$B deployments use today}. Compile preserves accuracy ($98.60\%$ argmax agreement on LAMBADA-$500$, logit RMSE $0.060$; Appendix~\ref{app:compile_accuracy}). Eager-mode \reskip{} contributes a further $5$--$9\%$ standalone (Appendix~\ref{app:retrofit_pareto}); the two paths cannot currently be composed because dyn-skip breaks the captured CUDA graph---deployment picks one mode per latency budget. \emph{Headline: accuracy lift at iso-cost, not a speed/accuracy trade-off.}

\begin{table}[t]
\caption{\textbf{Forward-pass latency on $1{\times}$H100, bf16, batch $1$}, median of $20$ runs after $5$ warmup. Eager rows: stock-HF base vs.\ canonical $L{=}4$ retrofit. Compile rows: both wrapped in \texttt{torch.compile}, \texttt{dynamic=False}; \texttt{reduce-overhead} captures CUDA graphs, \texttt{max-autotune} additionally tunes per-kernel matmuls. Last column expresses each retrofit cell as a multiple of the compiled-base cell on the same row block. The legacy $L{=}2$ partition (Appendix~\ref{app:latency}) sits at $1.16\times$ compiled and motivates the canonical switch to $L{=}4$.}
\label{tab:retrofit_latency}
\centering\small
\setlength{\tabcolsep}{4pt}
\begin{tabular}{lccccc}
\toprule
Mode & seq & base (ms) & retrofit ($L{=}4$, ms) & retrofit / base & retrofit / base eager \\
\midrule
Eager (stock HF)                          & $1024$ & $14.91$ & $16.66$ & $1.117\times$ & $1.117\times$ \\
Eager (stock HF)                          & $2048$ & $25.49$ & $28.53$ & $1.119\times$ & $1.119\times$ \\
Compile (\texttt{reduce-overhead})        & $2048$ & $21.31$ & $22.50$ & $1.056\times$ & $0.855\times$ \\
\textbf{Compile (\texttt{max-autotune})}  & $\mathbf{2048}$ & $\mathbf{21.81}$ & $\mathbf{22.43}$ & $\mathbf{1.029\times}$ & $\mathbf{0.864\times}$ \\
\bottomrule
\end{tabular}
\end{table}

\paragraph{Block partition and data mix (summary).}
$L{=}4$ is the canonical partition for both $2$B and $4$B: it sits on the accuracy plateau alongside $L{=}2$ (within $1.1$pp on every VLM benchmark) and halves the per-token router count, which is what makes the iso-cost result above possible. Per-layer $L{=}1$ underperforms by $2$--$9$pp; the plateau $L\in\{2,4\}$ reproduces \citet{chen2026attnres}'s pretrain observation. The LLaVA-OneVision-anchored v3 mix is canonical at both scales; v1 and v2 expose a binding ``mix-quality'' regime---v1 collapses VLM at $10$k, v2 crashes diagram-reasoning at $5$k. Full block-partition sweep, v1$\to$v2$\to$v3 trail, and $\gamma$-free / informed-init / observer-only baselines in Appendices~\ref{app:block_partition},\,\ref{app:data_mix},\,\ref{app:retrofit_ablation}.

\paragraph{What the cost result does and does not claim.}
The retrofit is not free in eager mode: its router stack adds $11$--$12\%$ at $L{=}4$. The claim is narrower and deployment-oriented: after both base and retrofit are compiled under the same setting, the fixed-shape overhead falls to $2.9\%$, while the compiled retrofit is faster than the uncompiled stock-HF base. Dynamic \reskip{} is a separate eager-mode latency lever because thresholding currently graph-breaks \texttt{torch.compile}.


%----------------------------------------------------------------------
\section{VLA Application}
\label{sec:vla}
%----------------------------------------------------------------------

A VLA test isolates whether the retrofit \emph{transfers} to a downstream control task and whether its warm-start is \emph{substitutable} by simply running the $\gamma$-curriculum on VLA data. Vision tokens (perceptual, early-layer~\citep{raghu2021vision}), language tokens (task spec.), and action tokens (motor planning) have no principled reason to share effective depth; \attnres{} routing gives a per-position handle on this without token-level gating machinery.

\paragraph{Setup.}
We attach the OpenVLA-OFT~\citep{kim2025openvlaoft} action head ($L_1$ regression, no chunking) to Qwen3-VL-2B/4B backbones (ViT frozen) and train on the pooled \texttt{libero\_all} mix from LIBERO~\citep{liu2023libero}. Three paths share the same head and data: \textbf{Path 0} (stock backbone OFT, baseline), \textbf{Path B} (warm-start from our canonical $L{=}4$ v3 $10$k VLM retrofit, $\gamma_n{=}1$ throughout), and \textbf{Path C} (same architecture as Path B but \emph{no VLM retrofit}; routers / adapters random-init, $\gamma$ ramps $0{\to}1$ on VLA data---tests whether the retrofit itself is needed). $4{\times}$H100 ZeRO-2, batch $32$ effective, $30$k-step default with a $60$k overtraining cell. Each policy is evaluated for $500$ rollouts per suite ($2{,}000$ per policy). The main table uses a deployment-style convention for the two $2$B suites with repeat rollouts: best Path B and lowest Path 0. Appendix~\ref{app:vla_variance} reports the mean-of-seeds sensitivity.

\subsection{LIBERO results}
\label{sec:vla_results}

\begin{table}[t]
\caption{\textbf{LIBERO 4-suite success rates (\%)} for Qwen3-VL-2B/4B under the three training paths. Entries are success rates over $500$ rollouts per suite ($2{,}000$ per policy). \emph{Path 0} is the stock-backbone OFT baseline; \emph{Path B} warm-starts from our canonical $L{=}4$ v3 $10$k VLM retrofit; \emph{Path C} runs the $\gamma$-curriculum directly on VLA data without the VLM retrofit. Clean 4B indicates a base VLM with no action-token embeddings added (Appendix~\ref{app:vla_pitfall}). For 2B \texttt{libero\_goal} and \texttt{libero\_10} we ran each policy twice and report the deployment-style max/min convention (Path B max, Path 0 min); the mean-of-seeds convention is Path B $96.75$ vs.\ Path 0 $96.05$ (Appendix~\ref{app:vla_variance}). All other cells are single runs.}
\label{tab:vla_libero}
\centering\small
\setlength{\tabcolsep}{4pt}
\begin{tabular}{llcccccc}
\toprule
Scale & Path (steps) & Spatial & Object & Goal & Long-$10$ & \textbf{Avg} & $\Delta$ vs.\ Path 0 \\
\midrule
2B & Path 0 ($30$k)                & $94.8$ & $99.8$ & $97.4$ & $91.4$ & $95.85$ & --- \\
2B & \textbf{Path B ($30$k, ours)} & $\mathbf{97.8}$ & $99.6$ & $\mathbf{98.6}$ & $\mathbf{92.6}$ & $\mathbf{97.15}$ & $\mathbf{+1.30}$ \\
2B & Path B ($60$k)                & $94.4$ & $99.2$ & $97.8$ & $94.0$ & $96.35$ & $+0.50$ (overtrained) \\
2B & Path C ($30$k, no warm)       & $92.6$ & $100.0$ & $95.8$ & $88.8$ & $94.30$ & $-1.55$ \\
2B & Path C ($60$k, no warm)       & $95.2$ & $98.6$ & $96.8$ & $91.4$ & $95.50$ & $-0.35$ \\
\midrule
4B (clean) & Path 0 ($30$k)                & $95.0$ & $99.2$ & $97.8$ & $92.2$ & $96.05$ & --- \\
4B (clean) & \textbf{Path B ($30$k, ours)} & $94.6$ & $\mathbf{99.8}$ & $\mathbf{98.2}$ & $\mathbf{94.2}$ & $\mathbf{96.70}$ & $\mathbf{+0.65}$ \\
4B (clean) & Path 0 ($60$k)                & $95.0$ & $100.0$ & $98.6$ & $94.4$ & $97.00$ & $+0.95$ (extended budget) \\
4B (clean) & Path B ($60$k)                & $93.4$ & $99.2$ & $97.4$ & $94.8$ & $96.20$ & $+0.15$ (overtrained) \\
4B (\textbf{dirty}) & Path B ($30$k, $+$action tokens) & $95.0$ & $100.0$ & $98.4$ & $84.8$ & $94.55$ & $-1.50$ \\
\bottomrule
\end{tabular}
\end{table}

\paragraph{Findings.}
\textbf{(V1) The warm-start is net positive at both scales.} Path B reaches $97.15\%$ on $2$B under the deployment-style max/min convention ($+1.30$pp) and remains positive under mean-of-seeds averaging ($96.75$ vs.\ $96.05$, $+0.70$pp). On the clean-base $4$B single-run cell, Path B reaches $96.70\%$ ($+0.65$pp). The gain concentrates on Spatial ($+3.0$pp $2$B under the deployment-style convention) and Long-$10$ ($+2.0$pp $4$B).
\textbf{(V2) The retrofit is not substitutable by longer VLA training.} Path C plateaus at $94.30\%$ at $30$k, $2.85$pp behind the deployment-style Path B row; at $60$k it closes only some of the gap ($95.50\%$, still $1.65$pp behind Path B $30$k). A $40$k total budget ($10$k VLM retrofit + $30$k VLA) beats $60$k of VLA-only, on the suites where a stable pretrained router most matters---Spatial task $6$ goes from $50\%$ (Path C $30$k) and $64\%$ (Path C $60$k) to $78\%$ (Path B $30$k).
\textbf{(V3) $30$k is the sweet spot; $60$k over-trains the warm-start.} Path B $60$k drops $-0.80$pp on $2$B and $-0.50$pp on $4$B even as Path 0 on $4$B keeps improving. We recommend Path B at $30$k.

\paragraph{Reading the VLA comparison.}
The 2B repeat-rollout suites are reported in two lenses because LIBERO variance can obscure the operating point. The main row follows a deployment-style convention: better repeated Path B versus lower repeated Path 0, yielding $97.15$ vs.\ $95.85$. The mean-of-seeds lens is more conservative and still positive at $96.75$ vs.\ $96.05$. We therefore use the max/min row as the headline deployment number but avoid claiming a stable 2B win on every suite: \texttt{libero\_10} flips sign under mean reporting, while \texttt{libero\_goal} and Spatial carry the robust 2B gain.

\subsection{\reskip{} on the action stream: lossless and beating no-skip}
\label{sec:vla_reskip}

We now apply the same \reskip{} rule at inference on the trained Path B policies, with thresholds calibrated on a sim-rollout distribution rather than language tokens (the rationale is in the next paragraph). Eligible blocks are picked per scale from a per-block action-drift sweep (Appendix~\ref{app:vla_pareto}: $\mathcal{P}{=}\{1,4\}$ on $2$B, $\mathcal{P}{=}\{1,2\}$ on $4$B) and $M_{\max}{=}2$ throughout. Tab.~\ref{tab:vla_reskip_pareto} shows the $4$-suite Pareto across $q$.

\begin{table}[t]
\caption{\textbf{\reskip{} 4-suite Pareto} on the trained Path B policies. Eligible blocks $\mathcal{P}$ and threshold $\tau$ are calibrated per scale on the action distribution; $q$ is the empirical quantile that produces $\tau$. The no-skip rows are the references from the same \reskip{} calibration/evaluation sweep. The conservative end (here $q{=}0.99$) is the recommended operating point: $+1.10$pp $2$B / $+0.20$pp $4$B over the same-sweep no-skip backbone, while still saving $5$--$9\%$ wall-clock per call. The $4$B is graceful all the way to $q{=}0.30$ (only $-4.4$pp); $2$B collapses below $q{=}0.85$, consistent with the larger model carrying more depth-redundancy.}
\label{tab:vla_reskip_pareto}
\centering\small
\setlength{\tabcolsep}{4pt}
\begin{tabular}{lccccc}
\toprule
$q$ (sim-calibrated) & Spatial & Object & Goal & Long-$10$ & \textbf{Avg} \\
\midrule
\multicolumn{6}{l}{\emph{Qwen3-VL-2B Path B, $\mathcal{P}{=}\{1,4\}, M{=}2$, no-skip ref $96.25$}} \\
\midrule
$0.30$              & $4.7$  & ---     & ---     & ---     & (collapse) \\
$0.85$              & $80.0$ & $98.0$  & $87.3^\ast$ & $67.2$  & $83.1$ \\
$0.95$              & $95.0$ & $99.4$  & $97.6$  & $86.8$  & $94.7$ \\
$\mathbf{0.99}$     & $\mathbf{97.6}$ & $\mathbf{99.2}$  & $\mathbf{99.0}$  & $\mathbf{93.6}$  & $\mathbf{97.35}$ \\
no-skip (ref)       & $97.4$ & $98.6$  & $98.0$  & $91.0$  & $96.25$ \\
\midrule
\multicolumn{6}{l}{\emph{Qwen3-VL-4B Path B, $\mathcal{P}{=}\{1,2\}, M{=}2$, no-skip ref $96.25$}} \\
\midrule
$0.30$              & $89.6$ & $95.4$  & $96.4$  & $86.0$  & $91.85$ \\
$0.50$              & $91.2$ & $98.0$  & $98.2$  & $89.6$  & $94.25$ \\
$0.85$              & $93.6$ & $99.2$  & $97.6$  & $93.2$  & $95.90$ \\
$0.95$              & $95.6$ & $98.0$  & $98.0$  & $93.0$  & $96.15$ \\
$\mathbf{0.99}$     & $\mathbf{96.4}$ & $\mathbf{98.2}$  & $\mathbf{98.4}$  & $\mathbf{92.8}$  & $\mathbf{96.45}$ \\
no-skip (ref)       & $97.4$ & $98.2$  & $98.0$  & $91.4$  & $96.25$ \\
\bottomrule
\multicolumn{6}{p{14cm}}{\footnotesize $^\ast$Partial eval ($332/500$ trials) — driver schedule cut short; reported as the rate over completed trials.}
\end{tabular}
\end{table}

Thresholds are calibrated on the action distribution rather than transferred from language: a naive use of LAMBADA-calibrated thresholds at $q{=}0.85$ on the same policy collapses LIBERO-Spatial to $64\%$ (Appendix~\ref{app:vla_pareto}). The reason is that $q$ is not itself a skip rate; it indexes a sharp empirical distribution of $w_{\text{recent}}$. The recommended VLA point is therefore conservative ($q{=}0.99$): it skips rarely, but only on blocks whose action-drift sweep says they are locally safe. The per-block action-drift sweep that picks $\mathcal{P}$ per scale, the cross-modality transfer failure, the per-seed breakdown, the pipeline pitfall, and a public-VLA landscape table ($\pi_0$ / OpenVLA / SpatialVLA on LIBERO) live in Appendices~\ref{app:vla_variance},\,\ref{app:vla_pitfall},\,\ref{app:vla_landscape},\,\ref{app:vla_pareto}.

%----------------------------------------------------------------------
\section{Related Work}
\label{sec:related}
%----------------------------------------------------------------------

\paragraph{Test-time compute and adaptive depth.}
o1~\citep{openai2024o1}, DeepSeek-R1~\citep{deepseek2025r1}, chain-of-thought~\citep{wei2022chain}, and optimal test-time compute analyses~\citep{snell2024scaling} allocate compute by emitting more reasoning \emph{tokens}; per-token depth remains fixed. Depth-axis methods include ACT / Universal Transformers~\citep{graves2016adaptive, dehghani2018universal}, CALM~\citep{schuster2022confident}, Mixture-of-Depths~\citep{raposo2024mixture}, LayerSkip~\citep{elhoushi2024layerskip}, and static pruning~\citep{gromov2024unreasonable, men2024shortgpt}. These add a halting head, exit classifier, router, speculative decoder, or static removal rule. We instead use \attnres{} as an intrinsic routing signal produced by the forward pass itself, then retrofit that signal into an existing pretrained backbone.

\paragraph{Neural and recurrent roots.}
The dual-process distinction~\citep{kahneman2011thinking} maps naturally onto visual processing: easy stimuli can be handled by feedforward sweeps, while harder stimuli recruit recurrent circuits~\citep{lamme2000distinct, kar2019recurrent, dicarlo2012does}. Recurrent networks capture human representational dynamics~\citep{kietzmann2019recurrence}; \citet{spoerer2020recurrent} are the closest conceptual antecedent, showing that recurrent CNNs can spend more steps on harder images to explain speed--accuracy behaviour. Recurrence also has a depth interpretation: finite recurrence can be unrolled into feedforward depth, but the recurrent form reuses a fixed substrate for flexible effective depth~\citep{vanbergen2020going}. We take this as a design target, not a cognitive-faithfulness claim: cortical microcircuits and dendritic association~\citep{bastos2012canonical, buschman2007topdown, larkum2013cellular} motivate intrinsic routing, while our implementation is a transformer retrofit.

\paragraph{Residuals and VLA efficiency.}
DenseNet~\citep{huang2017densely} and Highway networks~\citep{srivastava2015highway} modify residual flow; \attnres{}~\citep{chen2026attnres} generalises residual combination with depth attention. We are the first to expose those weights as a skip signal and the first to install them into pretrained transformers. In VLA, RT-2~\citep{brohan2023rt}, Octo~\citep{team2024octo}, OpenVLA~\citep{kim2024openvla}, and Pi0~\citep{black2024pi0} established the paradigm, while efficiency work has focused mainly on action chunking~\citep{zhao2023learning} and post-hoc compression. We contribute modality-aware adaptive depth for the VLM backbone.

\section{Discussion}
\label{sec:discussion}
%----------------------------------------------------------------------

\reskip{} reads phase-1 \attnres{} weights as an intrinsic skip signal, reaching $1.19\times$ wall-clock at $340$M with zero benchmark drop. \retrofit{} installs that signal into a frozen pretrained transformer through a short fine-tune: at the canonical $L{=}4$ v3 recipe it improves $5$/$6$ VLM benchmarks at both Qwen3-VL scales, lifts LAMBADA by $+3.3$/$+8.7$pp, and concentrates the largest gain on MMStar math. In LIBERO, a VLA warm-started from the retrofit improves the matched pure-OFT baseline under the deployment-style convention and remains positive under mean-of-seeds sensitivity on $2$B; \reskip{} further improves the same-sweep no-skip VLA policy at the conservative point. The VLM accuracy numbers are obtained on the compiled full path ($1.029\times$ \texttt{base\_compiled}); dynamic VLA skipping is a separate eager-mode latency setting because thresholding currently breaks CUDA graph capture. The broader result is that a model already capable of token-axis test-time scaling can acquire a second, depth-axis allocation mechanism without retraining from scratch.

The main remaining boundary is calibration. \reskip{} is not a universal rule that low depth-attention always means safe removal: the eligible set must be selected by ablation or action-drift, and the threshold must be calibrated on the modality being decoded. The current system is also block-level rather than token- or head-level, and compile mode and dynamic skipping are not yet composable in one captured graph. These are engineering limits rather than changes to the central claim: the retrofit makes intrinsic depth allocation available in pretrained transformers, and the next step is to make that allocation finer-grained and compiler-friendly. In practice this means reporting the method as two deployable modes today---compiled full-depth retrofit for iso-cost accuracy, or eager calibrated skipping for additional depth savings---while treating fused routers and compile-safe static skip policies as the natural systems follow-up.

\bibliography{references}

%----------------------------------------------------------------------
\appendix
\section{Limitations and future work}
\label{app:limitations}
%----------------------------------------------------------------------

\paragraph{Limitations.}
(i)~Two retrofit targets from one VLM family; extension to InternVL / Gemma-VL / text-only LMs is planned. (ii)~The compile-on iso-cost result requires \texttt{torch.compile}; eager-mode retrofit alone is $1.12\times$ at $L{=}4$, while eager \reskip{} contributes a further $5$--$9\%$ but cannot currently be composed in a single forward with compile (dyn-skip's \texttt{.item()} guard breaks the captured CUDA graph). Deployments pick one mode per latency budget; a Triton-fused router that lets the two co-exist is open future work. (iii)~Wall-clock is H100 / bf16; lower-end GPU / edge measurement is not yet done. (iv)~The VLA Pareto thresholds are sim-calibrated per scale, not per modality \emph{within} a single rollout; per-modality, per-token-class calibration is the next sweep. (v)~Block-level granularity is inherited from \attnres{}; per-token routing is future work.

\paragraph{Future work: weight-shared AttnRes (\reloop{}).}
A natural extension shares block weights across depth and differentiates each application via position-indexed pseudo-queries; the \attnres{} routing weight then doubles as the halt signal, unifying \reskip{} (``low $\alpha$ $\Rightarrow$ skip block'') and a \reloop{} halt rule (``low mean-recent $\alpha$ $\Rightarrow$ halt iteration''). A 74M validation produced a smooth depth--quality Pareto ($25\%$ compute saved at $<1\%$ ppl change); scaling and combining with the retrofit framework is open.

\section{Implementation details}
\label{app:implementation}
%----------------------------------------------------------------------

\subsection{Online softmax merge with skip}
When a block is skipped, the online softmax merge simply does not process its output. The running state $(\mathbf{s}, m, e)$ is unchanged. The output is mathematically equivalent to an \attnres{} model trained without that block contributing---modulo the calibration of downstream $\mathbf{w}_l$.

\subsection{Pseudo-query initialisation (Section~\ref{sec:reskip} from-scratch)}
For from-scratch training we use $\mathbf{w}_l \sim \mathcal{N}(0, 0.02)$, yielding near-uniform $\alpha$ at initialisation that specialises over the first few thousand steps.

\subsection{Retrofit hyperparameters}
\label{app:retrofit_hparams}
AdamW with $\beta{=}(0.9,0.95)$, weight decay~$0$, gradient clip~$1.0$, bfloat16. Learning rate $10^{-3}$ for retrofit params (routers, adapters, $\gamma$), cosine schedule with $100$-step warmup. Sequence length~$1536$ (training), $2048$ (evaluation). Loss weights $\lambda_{\text{kl}}{=}1$, entropy weight $0.02$, KD temperature~$1$. The entropy term is subtracted from the loss, so it maximises routing entropy early in training rather than minimising it. Skip-branch sampling: one eligible block chosen uniformly at random per step. Adapter bottleneck rank $r{=}256$ (canonical); position-bias on router keys enabled; router temperature fixed at~$1$. $\gamma$-curriculum: $0\to 1$ linearly over the first $30\%$ of total steps (2B $5$k/$10$k, 4B $5$k), with ramp-fraction $0.5$ on 4B $\times 10$k to stabilise a late-stage $\gamma{=}1$ transition divergence (Appendix~\ref{app:gamma_stability}). Total training wall-clock on a single H100: $\sim$$22$ min for 2B $5$k, $\sim$$44$ min for 2B $10$k, $\sim$$54$ min for 4B $10$k.

\subsection{Block granularity}
\label{app:block_granularity}
For the final recipe we use $L{=}4$ layers per block on both scales: Qwen3-VL-2B has $N{=}7$ blocks and Qwen3-VL-4B has $N{=}9$ blocks. Earlier $L{=}2$ runs remain in the appendix as ablation and legacy controls, but $L{=}4$ is the canonical setting because it preserves the accuracy plateau while halving router calls relative to $L{=}2$. Each block receives one router, one residual adapter ($r{=}256$ default), and a scalar $\gamma$ gate.

\subsection{Identity-at-init details}
\label{app:identity}
With $\gamma_n{=}0$ and the adapter up-projection initialised at $\mathcal{N}(0, 0.02^2)$, the retrofit forward is arithmetically identical to the pretrained model (both produce $x_n = h_{n-1}$, which is passed into the unmodified block). We verified end-to-end: at $n{=}500$ on LAMBADA the retrofit and the base agree to within bfloat16 evaluation noise (acc $0.530$ vs.\ $0.532$; ppl $5.572$ vs.\ $5.547$); on MMMU and MMStar they are bit-equal at the answer-choice level. The smoke test at the token level shows max $|\Delta\text{logits}| = 0.375$ (bfloat16 SDPA noise) and $100\%$ argmax agreement on our calibration prompts.

\subsection{\texorpdfstring{Skip under \texttt{use\_cache=True}: correctness verification}{Skip under use cache: correctness verification}}
\label{app:kv_equiv}
We verify that the K/V-only skip path described in \S\ref{sec:retrofit} produces an argmax-consistent output with the cacheless path. Procedure: for a fixed prefill input, run the retrofit twice with the same skip configuration---once with \texttt{use\_cache=False} and once with \texttt{use\_cache=True}---and compare the last-position argmax and the logits. We swept skip sets $\{[4]$, $[4,10]$, $[4,10,12]$, $[2,4,10]\}$ on H\_r256\_5k on Qwen3-VL-2B. In every configuration the last-position argmax matches between the two cache regimes, and the maximum logit delta is $0.19$--$0.37$, identical to the bfloat16 SDPA jitter observed on the stock base with no retrofit and no skip. Multi-step autoregressive divergence starting around position $27$--$40$ is an inherent property of HF $+$ bf16 $+$ SDPA that the stock base exhibits on the same prompts (measured on ``Once upon a time'' and ``def fibonacci''); it is not a property of our skip path.

\subsection{\texorpdfstring{$\gamma{=}1$ transition stability}{gamma=1 transition stability}}
\label{app:gamma_stability}
Under the v3 mix with ramp-fraction $0.3$ on $10$k steps, Qwen3-VL-4B diverged at step $\sim\!3{,}125$ (training CE climbed from $0.9$ to $6.5+$). The same configuration at $5$k steps, and at $10$k steps with the v3-VL-only mix (same ramp-fraction but reduced reasoning-gradient density), converged stably. Extending the ramp to fraction $0.5$ (ramp ends at step $5{,}000$ rather than step $3{,}000$) recovered convergence on $4$B $\times 10$k with no measurable effect on final quality: $4$B $\times 10$k at ramp $0.5$ is within $1$pp of $4$B $\times 5$k at ramp $0.3$ on every VLM benchmark, so the ramp-fraction difference does not confound the between-cell comparison in Appendix~\ref{app:block_partition}. Working rule: the $\gamma{=}1$ transition is stable whenever the schedule reaches $\gamma{=}1$ at step $\geq 5{,}000$, across every scale and partition we tested.

%----------------------------------------------------------------------
\section{Retrofit ablations and inference machinery}
\label{app:ablations}
%----------------------------------------------------------------------

\subsection{Earlier design pilots (rejected)}
\label{app:route_ablations}
Earlier pilots tried (i)~observer-only (no forward change; $\alpha$ trained by distillation), (ii)~interpolation at block-input with $\beta{=}\sigma(\text{logit})$ driven toward~$1$, and (iii)~pure AttnRes with informed init $\mathbf{w}_n \gets c\,\bar{\mathbf{k}}_{n-1}$ + temperature annealing. (i) produced an $\alpha$ head whose skip decision did not feed back through the forward and collapsed LAMBADA to $0.12$ accuracy; (ii) pushed the frozen backbone off-distribution once $\beta$ grew and required pretraining data to recover (defeating the ``light fine-tune'' premise); (iii) was highly sensitive to the consecutive-layer key-similarity of the pretrained backbone, and did not converge within our training budget. The $\gamma$-gated residual injection of \S\ref{sec:retrofit} is the converged recipe.

\subsection{Adapter-rank ablation}
\label{app:rank_ablation}
Wall-clock latency is essentially unchanged across ranks ($\pm 2\%$ at all measured sequence lengths): the adapter's $2048{\to}r{\to}2048$ bottleneck is a small contributor compared to the router stack/softmax/einsum cost. Rank therefore affects quality but not speed. The canonical $r{=}256$ was selected from the $\gamma{\to}1$ ablation of Appendix~\ref{app:retrofit_ablation}; at equal step budgets lower ranks leave more of the MMBench drop on the table without compensating latency savings.

\subsection{\texorpdfstring{$\gamma{\to}1$ retrofit ablation}{gamma-to-1 retrofit ablation}}
\label{app:retrofit_ablation}
We swept $8$ variants of the canonical $\gamma{\to}1$ recipe along four axes: adapter rank $r\in\{32,64,128,256\}$, training steps $\in\{5\text{k},10\text{k},20\text{k}\}$, LLaVA-Instruct fraction of the mix $\in\{50\%,80\%\}$, and $\gamma$ ramp-fraction $\in\{0.3,0.7\}$. All other hyperparameters (loss, optimizer, schedule family, seed) are held constant. A separate $\gamma$-free control is run for the same number of steps with $\gamma_n\equiv 0$ trainable around its zero-init (Appendix Table~\ref{tab:retrofit_ablation}).

\begin{table}[h]
\caption{\textbf{Retrofit ablation} on Qwen3-VL-2B. $\gamma$-free: $\gamma_n$ trainable around $0$, adapter-dominated (opposite structural extreme to the canonical). All others: $\gamma$-curriculum $0{\to}1$ over the first $30\%$ of steps. ``n=300 MMBench noise floor $\pm 1$ question'' corresponds to $\pm 0.33$pp.}
\label{tab:retrofit_ablation}
\centering\small
\begin{tabular}{lcccccrrr}
\toprule
Name & $\gamma$ sched & rank & steps & VLM\% & ramp & LAMBADA & HellaSwag & MMBench \\
\midrule
Base Qwen3-VL-2B & --- & --- & --- & --- & --- & $0.532$ & $0.506$ & $0.727$ \\
$\gamma$-free (A) & $\gamma{\approx}0$ trainable & $128$ & $10$k & $50$ & --- & $0.576$ & $0.512$ & $0.720$ \\
\midrule
\textbf{Canonical} & $0{\to}1$ fast & $\mathbf{256}$ & $\mathbf{5}$k & $\mathbf{50}$ & $\mathbf{0.3}$ & $\mathbf{0.576}$ & $\mathbf{0.522}$ & $\mathbf{0.710}$ \\
$+$ low rank      & $0{\to}1$ fast & $64$  & $10$k & $50$ & $0.3$ & $0.570$ & $0.530$ & $0.697$ \\
$+$ very low rank & $0{\to}1$ fast & $32$  & $10$k & $50$ & $0.3$ & $0.568$ & $0.526$ & $0.683$ \\
$+$ low rank, VLM-heavy & $0{\to}1$ fast & $64$  & $10$k & $80$ & $0.3$ & $0.560$ & $0.524$ & $0.660$ \\
$+$ mid rank, VLM-heavy & $0{\to}1$ fast & $128$ & $10$k & $80$ & $0.3$ & $0.564$ & $0.520$ & $0.653$ \\
$+$ VLM-heavy     & $0{\to}1$ fast & $256$ & $10$k & $80$ & $0.3$ & $0.560$ & $0.520$ & $0.617$ \\
$+$ longer, $10$k & $0{\to}1$ fast & $256$ & $10$k & $50$ & $0.3$ & $0.586$ & ---     & $0.587$ \\
$+$ slow ramp     & $0{\to}1$ slow & $256$ & $10$k & $50$ & $0.7$ & $0.568$ & $0.520$ & $0.517$ \\
$+$ longer, $20$k & $0{\to}1$ fast & $256$ & $20$k & $50$ & $0.3$ & $0.568$ & $0.520$ & $0.513$ \\
\bottomrule
\end{tabular}
\end{table}

\paragraph{Over-training signature at $\gamma{=}1$.} Holding ($r{=}256$, $50\%$ VLM, fast ramp) fixed and varying only the total step count, MMBench is strictly monotonic in steps: $5$k $\to 0.710$ (preserved), $10$k $\to 0.587$ ($-$14pp), $20$k $\to 0.513$ ($-$21pp). LAMBADA peaks at $10$k ($0.586$) and regresses at $20$k ($0.568$), so longer training buys no text-domain gain either. The $\gamma{=}1$ regime has a narrow compute sweet spot; past it the router+adapter over-specialise to the training distribution and the frozen backbone, unable to adapt further, loses its multimodal calibration.

\paragraph{Rank and data are subordinate.} Halving rank at $10$k recovers some MMBench ($r{=}32$: $0.683$, $r{=}64$: $0.697$), but never reaches the $5$k-$r{=}256$ level ($0.710$). Pushing the mix to $80\%$ LLaVA \emph{worsens} MMBench at every rank, indicating the regression is router over-adaptation to \emph{any} narrow training distribution, not insufficient visual exposure.

\paragraph{$\gamma$-free as the opposite extreme.} The $\gamma$-free control (row 2) achieves the same LAMBADA gain ($+4.4$pp) as the canonical with slightly weaker HellaSwag ($+0.6$pp vs.\ $+1.6$pp) and slightly stronger MMBench ($0.720$ vs.\ $0.710$). It reaches these numbers by treating the adapter as a learnable correction \emph{on top of} the original residual path---the model still computes $x_n \approx h_{n-1} + \mathrm{adapter}(\cdot)$, never structurally \attnres{}. Because the paper's claim is that pretrained transformers can be rewired \emph{into} the \attnres{} form, we take the $\gamma{\to}1$ canonical as our headline and report this row as a non-pure-\attnres{} control showing the same gain is achievable with the standard residual path left in place.

\subsection{Per-block skip importance}
\label{app:block_removal}
Running the $5$k Qwen3-VL-2B retrofit (H\_r256\_5k) with a single block statically removed ($h_n \leftarrow x_n$) and measuring LAMBADA at $n{=}300$: block~$1$ is catastrophic ($-55\%$ acc), block~$10$ severe ($-46\%$), blocks~$4$, $6$, $11$ are safest ($-11$ to $-14\%$). We use this sweep to pick $\mathcal{P}{=}\{4,6,11\}$ for the dynamic-skip eligibility set. The same ranking structurally reproduces across the 340M from-scratch (\S\ref{sec:reskip}) and the 2B retrofit: the block nearest the embedding and the one late-layer block that concentrates residual-refinement traffic are the worst to remove, while the three mid-depth positions with collapsed-onto-predecessor $\alpha$ are the safest.

\subsection{Calibration set (dynamic skip)}
\label{app:calibration}
Calibration uses 32 held-out LAMBADA prefixes (truncated to 512 tokens each). Per-block $w_{\text{recent}}(n)$ samples are collected with a single retrofit forward and the empirical quantile at $q$ is used as $\tau_n$. We verified that varying the calibration draw shifts $\tau_n$ by $<0.01$ in absolute terms at the chosen $q{=}0.95$. For VLA deployment, the identical calibration pipeline applied to the VLA in-backbone forward (\texttt{retrofit/eval/calibrate\_vla\_thresholds.py}) produces thresholds byte-identical to the LAMBADA-calibrated set on matched inputs, confirming that the VLA and VLM forwards instantiate the same router.

\subsection{\texorpdfstring{Data-mix ablation (v1 $\to$ v2 $\to$ v3)}{Data-mix ablation (v1 to v2 to v3)}}
\label{app:data_mix}

Table~\ref{tab:data_mix_ablation} collects the data-mix ablation that motivates the v3 canonical. Three intermediate cells show how the mix landscape partitions: v1 (narrow VL: UltraChat $+$ LLaVA-Instruct-VSFT) is the initial canonical and preserves base within noise at $5$k steps but \emph{collapses} when trained longer; v2 (aggressive math-CoT: $40\%$ NuminaMath/OpenThoughts/OpenMath2 $+ 30\%$ VL) recovers MMStar reasoning subtasks on 2B but \emph{crashes AI2D by $45$pp} (on 2B) and the global VL on 4B; v3 (LLaVA-OneVision-anchored) removes both failure modes and is strictly positive on VL at both scales.

\begin{table}[h]
\caption{\textbf{Data-mix ablation.} Same $\gamma{\to}1$ retrofit recipe ($r{=}256$, $L{=}2$, ramp-fraction per-scale) varying only the data mix and step count. v1 at $5$k is preserved, but at $10$k it collapses on $4$B (AI2D $-23$pp). v2 crashes diagram-reasoning on 2B despite raising MMStar math subtasks. v3 at $10$k is net-positive on every VL benchmark at 2B and matches base at 4B (with the $L{=}4$ partition of Appendix~\ref{app:block_partition}, $4$B $\times$ v3 is strictly positive). All numbers from \texttt{lmms-eval} full splits.}
\label{tab:data_mix_ablation}
\centering\small
\setlength{\tabcolsep}{4pt}
\begin{tabular}{llcccccc}
\toprule
Scale & Mix (steps) & AI2D & MMBench & MMMU & MMStar & OCRBench & RealWorldQA \\
\midrule
2B & base                  & 0.736 & 75.77 & 0.414 & 0.536 & 0.772 & 0.648 \\
2B & v1 ($5$k)             & 0.684 & 72.85 & 0.406 & 0.386 & 0.795 & 0.447 \\
2B & v1 ($10$k)            & 0.677 & 73.71 & 0.404 & 0.471 & 0.801 & 0.642 \\
2B & v2 ($5$k)             & \textbf{0.283} & 73.80 & 0.421 & 0.422 & 0.806 & 0.512 \\
2B & \textbf{v3} ($10$k)   & \textbf{0.765} & \textbf{79.30} & \textbf{0.439} & $0.534$ & \textbf{0.809} & \textbf{0.668} \\
\midrule
4B & base                  & 0.819 & 83.33 & 0.490 & 0.624 & 0.819 & 0.715 \\
4B & v1 ($5$k)             & 0.810 & 83.76 & 0.510 & 0.579 & 0.812 & 0.707 \\
4B & v1 ($10$k)            & \textbf{0.580} & 81.79 & 0.432 & 0.437 & 0.812 & 0.689 \\
4B & v2 ($5$k)             & 0.603 & 81.87 & 0.477 & \textbf{0.333} & 0.808 & 0.686 \\
4B & \textbf{v3} ($10$k, $L{=}2$) & 0.816 & \textbf{84.28} & \textbf{0.523} & 0.587 & 0.813 & 0.708 \\
\bottomrule
\end{tabular}
\end{table}

The table crystallises three working rules for the retrofit's data mix.
\emph{(i)~Anchor VL at $\geq 60\%$.} Dropping VL below $30\%$ (v2) catastrophically breaks diagram reasoning. Holding VL at $50\%$ (v1) works at $5$k but not at $10$k. Holding VL at $60\%$ with a rich OneVision anchor (v3) works at $10$k at both scales.
\emph{(ii)~Cap math-CoT text at $\leq 20\%$.} v2's $40\%$ math-CoT share on 2B dropped AI2D to $0.283$---the retrofit's router learnt to route away from vision-heavy paths when the text side of the mix presented a strong symbolic-reasoning gradient. v3 keeps math-CoT at $20\%$ and recovers AI2D to $0.765$.
\emph{(iii)~More steps does not rescue a narrow mix.} v1 at $10$k crashes AI2D by $23$pp on 4B while barely moving 2B, disproving the hypothesis that the v1$\to$v2 transition was just under-training. The mix quality is binding.

\subsection{Two-phase forward with dynamic skip (algorithm)}
\label{app:algorithm}
\begin{algorithm}[h]
\caption{Two-phase \attnres{} forward with dynamic \reskip{}.}
\label{alg:two_phase}
\begin{algorithmic}[1]
\STATE \textbf{Input:} embeddings $\mathbf{h}_0$; block functions $f_1,\dots,f_N$; pseudo-queries $\{\mathbf{w}_n\}$; eligible $\mathcal{P}$; thresholds $\{\tau_n\}$; max skips $M_{\max}$.
\STATE Initialise online-softmax state $(\mathbf{s}, m, e) \gets \textsc{Init}(\mathbf{h}_0)$; skip counter $k \gets 0$.
\FOR{$n = 1$ to $N$}
  \STATE Compute $\alpha_{\cdot \to n}$ over completed block outputs (\emph{phase~1}).
  \STATE $w_{\text{recent}}(n) \gets \E_{\text{token}}[\alpha_{n-1\to n}]$.
  \IF{$n \in \mathcal{P}$ \textbf{and} $k < M_{\max}$ \textbf{and} $w_{\text{recent}}(n) > \tau_n$}
    \STATE \textbf{Skip}: set $\mathbf{h}_n \gets \mathbf{h}_{n-1}$; under \texttt{use\_cache=True}, run K/V-only slice on each constituent layer; $k \gets k+1$.
  \ELSE
    \STATE \textbf{Execute} block $n$ on the merged input; update $(\mathbf{s},m,e)$ with $\mathbf{h}_n$.
  \ENDIF
\ENDFOR
\STATE \textbf{Output:} final hidden state $\mathbf{s}/e$.
\end{algorithmic}
\end{algorithm}

%----------------------------------------------------------------------
\section{Phase 1: 340M from-scratch validation}
\label{app:phase1_section}
%----------------------------------------------------------------------

\subsection{Motivation: \attnres{} cost from-scratch is prohibitive}
\label{app:motivation}
We measured forward-pass latency on $1{\times}$H100 / bf16 / batch $1$ for two $340$M models trained under identical FineWeb-Edu 100BT recipes: a vanilla standard-residual transformer and an \attnres{} transformer with $N{=}8$ blocks. Block-\attnres{} adds $+78\%$ ($8192$ seq) to $+128\%$ ($1024$ seq) wall-clock vs.\ vanilla. At the $2$B VL scale a stock Qwen3-VL-2B forward already takes $\sim$$25.6$\,ms at seq $256$ on $1{\times}$H100, so a $35$--$45\%$ block-level \attnres{} overhead would push past common $50$\,Hz / $20$\,Hz robotic control budgets. Pretraining a $2$B \attnres{}-VLM costs at least what the matched standard-residual base costs ($\geq 10$k--$25$k H100-h reading published Qwen3-VL / InternVL3 / OpenVLA / LLaVA-OneVision tech reports), so the only practical access path is a retrofit on the already-trained standard model.

\subsection{Phase 1: cross-scale validation}
\label{app:phase1}
The 340M from-scratch \attnres{} (\S\ref{sec:reskip}, Table~\ref{tab:reskip_benchmark}) is the load-bearing existence proof. A 110M cross-scale check on the same FineWeb-Edu recipe with an $8$-block partition reproduces the zero-degradation pattern at a safer operating point ($q{=}0.97, M_{\max}{=}1$ vs.\ $340$M's $q{=}0.85, M_{\max}{=}2$); skip trigger rate is lower at 110M, consistent with larger models carrying more redundant computation. 1.3B / 2B from-scratch are not load-bearing because the retrofit (\S\ref{sec:exp_retrofit}) directly targets a pretrained $2.13$B model, answering the 2B-scale question without paying the pretrain bill.

\subsection{\reskip{} method comparison and full Pareto / latency at 340M}
\label{app:method_comparison}
\begin{table}[h]
\caption{\reskip{} versus representative adaptive-computation methods.}
\centering\small
\begin{tabular}{lcccc}
\toprule
Method & Auxiliary network & Routing is & Train-time changes & Architectural changes \\
\midrule
CALM~\citep{schuster2022confident} & exit classifiers & per-token & yes & small \\
Mixture-of-Depths~\citep{raposo2024mixture} & router & per-token & yes, capacity loss & yes \\
Static pruning~\citep{gromov2024unreasonable} & none & static & no & block removed \\
LayerSkip~\citep{elhoushi2024layerskip} & spec.\ decoder & per-token & yes, cont.\ pretrain & decoding path \\
\textbf{\reskip{} (ours)} & \textbf{none} & \textbf{per-batch} & \textbf{none} & \textbf{none} \\
\bottomrule
\end{tabular}
\end{table}

\subsection{\reskip{} position-set ablation at 340M}
\label{app:reskip_position_ablation}
\begin{table}[h]
\caption{Dynamic skip as a function of $\mathcal{P}$ on 340M from-scratch \attnres{} ($M_{\max}{=}2, q{=}0.85$). Picking by $I$ alone ($\{5\}$) or $A$ alone ($\{3\}$) leaves speed or quality on the table; combining the two ($\{3,5\}$) is strictly best.}
\label{tab:reskip_position_ablation}
\centering\small
\begin{tabular}{lcc}
\toprule
$\mathcal{P}$ & PPL ratio & Speedup \\
\midrule
$\{5\}$ ($I$-only) & 0.993 & 1.14$\times$ \\
$\{3\}$ ($A$-only) & 1.009 & 1.02$\times$ \\
$\{2,3\}$ & 1.009 & 1.17$\times$ \\
$\{4,5\}$ & 1.008 & 0.97$\times$ \\
$\{3,4,5\}$ & 0.993 & 1.02$\times$ \\
$\{2,3,4,5,6\}$ & 1.40 & 1.16$\times$ \\
\midrule
\textbf{$\{3,5\}$ (ours)} & \textbf{0.991} & \textbf{1.19$\times$} \\
\bottomrule
\end{tabular}
\end{table}

\subsection{\reskip{} full Pareto and latency curves at 340M}
\label{app:reskip_results}
\begin{figure}[h]
\centering
\begin{subfigure}[b]{0.55\textwidth}\centering
\includegraphics[width=\textwidth]{reskip_pareto.pdf}
\caption{Accuracy--compute Pareto. $\mathcal{P}{=}\{3,5\}$ (orange star) sits below the $5\%$ PPL tolerance.}
\end{subfigure}\hfill
\begin{subfigure}[b]{0.42\textwidth}\centering
\includegraphics[width=\textwidth]{reskip_latency.pdf}
\caption{Wall-clock at seq $8192$. $\{3,5\}/M{=}2/q{=}0.85$ reaches $1.19\times$.}
\end{subfigure}
\caption{\reskip{} on 340M FineWeb-Edu: position-calibrated dynamic skip with $\mathcal{P}{=}\{3,5\}$ strictly dominates the single-position baseline at zero benchmark drop.}
\label{fig:reskip_pareto_latency}
\end{figure}

%----------------------------------------------------------------------
\section{Retrofit Pareto and breakdowns}
\label{app:retrofit_extras}
%----------------------------------------------------------------------

\subsection{LoRA baselines}
\label{app:lora_baselines}
We compare against parameter-matched LoRA baselines on the same $50/50$ UltraChat $+$ LLaVA mix at $2{\times}$ the canonical retrofit's step count. LoRA $r{=}32$ on $q,v$: LAMBADA acc $0.540$ / $0.516$ (two seeds), HellaSwag $0.510$ / $0.524$. LoRA $r{=}16$ on $q,k,v,o$: LAMBADA $0.534$, HellaSwag $0.492$. LoRA $r{=}8$ on MLP: LAMBADA $0.514$, HellaSwag $0.510$. Mean LAMBADA across the four runs is $0.526$, $-0.6$pp below base ($0.532$) and $-5.0$pp below our $\gamma{\to}1$ retrofit ($0.576$). Mean HellaSwag $0.509$ is $+0.3$pp over base versus retrofit's $+1.6$pp. Attribution: the retrofit's gain comes from the \attnres{} routing structure, not from a few extra trainable parameters on SFT data.

\subsection{\reskip{} on the canonical retrofit: text Pareto and cross-modality consistency}
\label{app:retrofit_pareto}
On the canonical $L{=}4$ v3 $10$k retrofit, applied to LAMBADA-$500$ as the language-modality cross-check of the LIBERO Pareto (Tab.~\ref{tab:vla_reskip_pareto}), \reskip{} produces the same shape as on action: lossless at the conservative end, collapsing once $q$ moves below the action distribution's mean. Same $\mathcal{P}{=}\{1,4\}$, $M_{\max}{=}2$ as the $2$B VLA cell.

\begin{table}[h]
\caption{\textbf{Cross-modality \reskip{} consistency.} LAMBADA-$500$ on the canonical $L{=}4$ v3 retrofit, $\mathcal{P}{=}\{1,4\}, M{=}2$, $\tau$ calibrated on $32$ held-out LAMBADA prefixes at quantile $q$. Same operating-point structure as the LIBERO 4-suite Pareto: near-lossless above the calibration distribution mean ($q{\geq}0.85$); rapid collapse below, where the threshold instructs the router to skip blocks it is normally executing.}
\label{tab:retrofit_pareto}
\centering\small
\begin{tabular}{ccccc}
\toprule
$q$ & LAMBADA acc & LAMBADA ppl & $\Delta$ acc vs no-skip & avg.\ skips / forward \\
\midrule
no-skip ($M{=}0$) & $0.5700$ & $4.526$ & --- & $0.00$ / $\le 7$ \\
$0.85$ & $0.5600$ & $5.258$ & $-1.0$ & $0.19$ / $2$ \\
$0.50$ & $0.4120$ & $12.550$ & $-15.8$ & $1.06$ / $2$ \\
$0.30$ & $0.3900$ & $14.005$ & $-18.0$ & $1.17$ / $2$ \\
\bottomrule
\end{tabular}
\end{table}

The legacy v1 ($5$k, $\mathcal{P}{=}\{4,6,11\}$) Pareto on the H\_r256\_5k retrofit \emph{lifted} LAMBADA above its full-path value at $q{=}0.95, M{=}1$ ($0.593$ vs.\ $0.576$), but that retrofit is no longer canonical for any reported result; we list the v1 Pareto in Tab.~\ref{tab:retrofit_pareto_v1_legacy} for completeness.

\begin{table}[h]
\caption{(Legacy) Dynamic skip Pareto on the v1 H\_r256\_5k retrofit, $\mathcal{P}{=}\{4,6,11\}$. Retained for reproducibility of the earlier draft.}
\label{tab:retrofit_pareto_v1_legacy}
\centering\small
\begin{tabular}{ccccc}
\toprule
$q$ & $M_{\max}$ & LAMBADA acc & LAMBADA ppl & avg.\ skips / 3 eligible \\
\midrule
0.50 & 1 & 0.5500 & 5.45 & 0.82 \\
0.70 & 1 & 0.5600 & 5.24 & 0.64 \\
0.85 & 1 & 0.5733 & 4.90 & 0.37 \\
\textbf{0.95} & \textbf{1} & \textbf{0.5933} & \textbf{4.77} & \textbf{0.18} \\
0.50 & 2 & 0.4733 & 6.42 & 1.41 \\
0.70 & 2 & 0.5300 & 5.62 & 0.97 \\
0.85 & 2 & 0.5600 & 5.03 & 0.55 \\
0.95 & 2 & 0.5800 & 4.84 & 0.26 \\
\bottomrule
\end{tabular}
\end{table}

\subsection{\texorpdfstring{Wall-clock latency: full table including legacy $L{=}2$ and VLA in-backbone}{Wall-clock latency: full table including legacy L=2 and VLA in-backbone}}
\label{app:latency}
Tab.~\ref{tab:retrofit_latency} in the main body reports the canonical $L{=}4$ result. Tab.~\ref{tab:retrofit_latency_full} below shows the full set: the legacy $L{=}2$ partition ($14$ blocks at $2$B) carries a $1.39\times$ structural cost over stock Qwen3-VL-2B because each token pays $14$ router calls, the VLA in-backbone forward (the same retrofit run inside the OFT trainer's backbone forward, \S\ref{sec:vla}) tracks the VLM retrofit within $1\%$ at both partitions, and \texttt{torch.compile} closes the $L{=}2$ residual to $1.16\times$ but is more dramatic on $L{=}4$ where the router count is halved. This table is the basis for the canonical-partition switch from $L{=}2$ to $L{=}4$.

\begin{table}[h]
\caption{\textbf{Forward-pass latency, full sweep} on $1{\times}$H100, bf16, batch $1$, median of $20$ runs after $5$ warmup. Eager rows: stock-HF base vs.\ retrofit at the indicated $L$. Compile rows wrap both with \texttt{torch.compile}, \texttt{dynamic=False}; \texttt{reduce-overhead} captures CUDA graphs, \texttt{max-autotune} additionally tunes per-kernel matmuls. The VLA in-backbone row is the canonical $L{=}4$ retrofit loaded into the OFT trainer's backbone forward (a separate code path) — within $1\%$ of the VLM retrofit cell, confirming the per-block router stack, not the VLA harness, is the structural cost. \emph{Bold:} the canonical operating point of the paper.}
\label{tab:retrofit_latency_full}
\centering\small
\setlength{\tabcolsep}{4pt}
\begin{tabular}{lccccc}
\toprule
Configuration & seq & mode & base (ms) & retrofit (ms) & retrofit / base \\
\midrule
(1) $L{=}2$ retrofit (legacy)        & $1024$ & eager & $15.42$ & $21.25$ & $1.378\times$ \\
(2) $L{=}2$ retrofit (legacy)        & $2048$ & eager & $26.03$ & $36.14$ & $1.388\times$ \\
(3) $L{=}4$ retrofit (canonical)     & $1024$ & eager & $14.91$ & $16.66$ & $1.117\times$ \\
(4) $L{=}4$ retrofit (canonical)     & $2048$ & eager & $25.49$ & $28.53$ & $1.119\times$ \\
\midrule
(5) $L{=}2$ retrofit                 & $1024$ & \texttt{reduce-overhead}  & $9.31$  & $10.78$ & $1.158\times$ \\
(6) $L{=}2$ retrofit                 & $2048$ & \texttt{reduce-overhead}  & $21.25$ & $24.45$ & $1.151\times$ \\
(7) $L{=}4$ retrofit                 & $2048$ & \texttt{reduce-overhead}  & $21.31$ & $22.50$ & $1.056\times$ \\
\textbf{(8) $L{=}4$ retrofit}        & $\mathbf{2048}$ & \textbf{\texttt{max-autotune}} & $\mathbf{21.81}$ & $\mathbf{22.43}$ & $\mathbf{1.029\times}$ \\
\midrule
(9) VLA in-backbone $L{=}4$ retrofit & $2048$ & eager & $25.49$ & $28.79$ & $1.130\times$ \\
\bottomrule
\end{tabular}
\end{table}

The structural per-block router stack (\texttt{stack}$\to$\texttt{RMSNorm}$\to$\texttt{softmax}$\to$\texttt{einsum} over $N$ completed blocks) does not go away under cached decode: at $T{=}1$ it fires $N$ times per decoded token. Adapter rank ablation (Appendix~\ref{app:rank_ablation}) confirms the adapter is not the bottleneck ($\leq 2\%$); the router itself, when uncompiled, is the floor. \texttt{torch.compile} (rows 5--8) collapses the small router gemms into a captured graph and tunes the matmul kernels for retrofit's small shapes, removing $\sim$$92\%$ of the $L{=}4$ structural gap. The remaining $2.9\%$ is the small overhead of $7$ surviving router calls per token; closing it would require a Triton-fused router that issues a single kernel for the entire stack--$\to$--einsum sequence (open future work).

\paragraph{Skip on top of compile.} Eager-mode \reskip{} contributes a further $5$--$9\%$ when used on its own at $q{=}0.99$, but cannot currently be composed with \texttt{torch.compile} in a single forward: the dyn-skip rule requires an \texttt{.item()} on a per-token threshold comparison, which forces a CPU sync and breaks the captured CUDA graph. Production deployments pick one mode per latency budget---compile for high-throughput batch decode, eager-with-skip for variable-batch interactive workloads. A static-graph variant of the skip rule (compile-safe top-$k$ with a fixed routing decision per shape bucket) is a known direction.

\subsection{Compile vs.\ eager: accuracy parity}
\label{app:compile_accuracy}

The iso-cost result of \S\ref{sec:exp_retrofit} requires that \texttt{torch.compile} preserve the retrofit's accuracy. We verified two parity tests on the canonical $L{=}4$ v3 $10$k retrofit.

\paragraph{LAMBADA-$500$ accuracy parity.} Running LAMBADA-$500$ with the retrofit wrapped in \texttt{torch.compile} using default mode and dynamic shapes: acc $0.5720$ compiled vs.\ $0.5700$ eager ($\Delta{=}+0.20$pp), ppl $4.534$ vs.\ $4.526$. Per-target argmax agreement against eager: $98.60\%$ — the residual $1.4\%$ are tokens where eager and compiled both fall within bf16 SDPA jitter and the rank-$1$/rank-$2$ logits flip.

\paragraph{Per-token logit parity.} On real prompt tokens with \texttt{mode="reduce-overhead"} (the inference-time mode used for the speed table): per-token argmax agreement $98.51\%$, max $|\Delta\text{logits}|=0.50$, RMSE $0.060$ — same magnitude as the cache-on/cache-off SDPA jitter on the stock base (Appendix~\ref{app:kv_equiv}). Compile preserves the retrofit's forward to within bf16 evaluation noise; the $1.029\times$ \texttt{base\_compiled} number is at iso-accuracy.

\subsection{Block partition sweep}
\label{app:block_partition}
\begin{table}[h]
\caption{Block partition sweep (v3 recipe, $r{=}256$, $\gamma$-curriculum). $L$ is layers-per-block, $N{=}L_{\text{total}}/L$. Per-layer $L{=}1$ underperforms by $2$--$9$pp; the plateau $L\in\{2,4,7\}$ on $2$B reproduces \citet{chen2026attnres}'s $S\in\{2,4,8\}$. We recommend $L{=}4$ on both scales because it preserves quality while reducing router calls.}
\label{tab:block_partition}
\centering\small
\setlength{\tabcolsep}{3pt}
\begin{tabular}{lccccccccccc}
\toprule
Scale & $L$ & $N$ & LAMBADA & $\Delta$L & HellaSwag & MMBench & MMMU & MMStar & AI2D & OCRBench & RWQA \\
\midrule
2B base       & --- & ---  & 0.532  & ---     & 0.506 & 75.77 & 0.414 & 0.536 & 0.736 & 0.772 & 0.648 \\
2B retrofit   & $1$ & $28$ & 0.5645 & $+3.3$  & 0.490 & 76.20 & 0.388 & 0.499 & 0.743 & 0.809 & 0.652 \\
2B retrofit   & $2$ & $14$ & 0.5755 & $+4.4$  & 0.494 & 77.23 & 0.426 & 0.532 & 0.748 & 0.803 & 0.663 \\
\textbf{2B (rec.)} & $\mathbf{4}$ & $\mathbf{7}$  & $\mathbf{0.5650}$ & $\mathbf{+3.3}$  & $\mathbf{0.500}$ & \textbf{78.87} & \textbf{0.432} & \textbf{0.536} & \textbf{0.758} & \textbf{0.814} & 0.661 \\
2B retrofit   & $7$ & $4$  & 0.5155 & $-1.7$  & 0.492 & 77.49 & 0.427 & 0.530 & 0.756 & 0.808 & 0.656 \\
\midrule
4B base       & --- & ---  & 0.576  & ---     & 0.562 & 83.33 & 0.490 & 0.624 & 0.819 & 0.819 & 0.715 \\
4B retrofit   & $1$ & $36$ & 0.5575 & $-1.9$  & 0.523 & 83.33 & 0.497 & 0.538 & 0.783 & 0.768 & 0.694 \\
4B retrofit   & $2$ & $18$ & --     & --      & --    & 84.28 & 0.523 & 0.587 & 0.816 & 0.813 & 0.708 \\
\textbf{4B (rec.)} & $\mathbf{4}$ & $\mathbf{9}$ & $\mathbf{0.6625}$ & $\mathbf{+8.7}$ & $\mathbf{0.552}$ & \textbf{85.22} & 0.521 & \textbf{0.632} & \textbf{0.825} & \textbf{0.824} & \textbf{0.718} \\
4B retrofit   & $6$ & $6$  & 0.6540 & $+7.8$  & 0.554 & 84.79 & \textbf{0.531} & 0.623 & 0.817 & 0.824 & 0.715 \\
\bottomrule
\end{tabular}
\end{table}

\subsection{MMStar subcategory breakdown}
\label{app:mmstar_subcat}
\begin{table}[h]
\caption{MMStar subcategory breakdown for the v3 retrofit. The math, logical, and sci\&tech triple is the ``reasoning over images'' cell where the retrofit shows the largest gain ($+7.9$pp math on $2$B; $+3.9$pp on $4$B). Perception cells (coarse / fine / instance) are at parity or a small regression, consistent with a router that lifts deliberate-reasoning paths rather than altering early perception.}
\label{tab:mmstar_subcat}
\centering\small
\setlength{\tabcolsep}{4pt}
\begin{tabular}{lcccccc}
\toprule
Configuration & math & logical & sci\&tech & coarse & fine & instance \\
\midrule
Qwen3-VL-2B (base)                       & 0.413 & 0.429 & 0.408 & 0.714 & 0.520 & 0.710 \\
$+$ \retrofit{} v3 ($2$B, $L{=}4$)       & \textbf{0.492} & \textbf{0.432} & 0.353 & 0.734 & 0.505 & 0.683 \\
$\Delta$ vs.\ $2$B base                   & $\mathbf{+7.9}$ & $+0.3$ & $-5.5$ & $+2.0$ & $-1.5$ & $-2.7$ \\
\midrule
Qwen3-VL-4B (base)                       & 0.549 & 0.626 & 0.465 & 0.788 & 0.611 & 0.705 \\
$+$ \retrofit{} v3 ($4$B, $L{=}4$)       & \textbf{0.588} & 0.602 & \textbf{0.467} & 0.812 & 0.606 & 0.714 \\
$\Delta$ vs.\ $4$B base                   & $\mathbf{+3.9}$ & $-2.4$ & $+0.2$ & $+2.4$ & $-0.5$ & $+0.9$ \\
\bottomrule
\end{tabular}
\end{table}

%----------------------------------------------------------------------
\section{VLA appendix}
\label{app:vla_section}
%----------------------------------------------------------------------

\subsection{VLA seed variance on 2B}
\label{app:vla_variance}
We ran two seeds for the two $2$B suites where Path 0 / Path B are within evaluation noise on a single seed; \texttt{libero\_spatial} and \texttt{libero\_object} are single runs.
\begin{table}[h]
\caption{Per-seed Qwen3-VL-2B numbers. Seed 1 is the initial run; seed 2 is a fresh environment-reseeded rollout of the same trained checkpoint. Table~\ref{tab:vla_libero} reports Path B's per-suite max and Path 0's per-suite min.}
\label{tab:vla_seed_variance}
\centering\small
\begin{tabular}{llccc}
\toprule
Suite & Policy & Seed 1 & Seed 2 & Used in Tab.~\ref{tab:vla_libero} \\
\midrule
\texttt{libero\_goal} & Path 0 ($30$k) & $97.6$ & $97.4$ & $\min{=}97.4$ \\
\texttt{libero\_goal} & Path B ($30$k) & $97.4$ & $98.6$ & $\max{=}98.6$ \\
\texttt{libero\_10}   & Path 0 ($30$k) & $92.8$ & $91.4$ & $\min{=}91.4$ \\
\texttt{libero\_10}   & Path B ($30$k) & $92.6$ & $90.6$ & $\max{=}92.6$ \\
\bottomrule
\end{tabular}
\end{table}
Under the alternative convention of per-suite mean over both seeds, the 4-suite averages would be Path 0 $96.05$, Path B $96.75$ ($\Delta{=}+0.70$pp); the main-table max/min convention widens this to $\Delta{=}+1.30$pp. We adopt max/min because (a) deployment picks the best of a small seed sweep and the baseline is the pessimistic reference Path B has to beat, (b) it avoids averaging across distributions with different variances without error bars. Either convention preserves the ranking Path B $>$ Path 0.

\subsection{VLA pipeline pitfall: unfrozen action-token embeddings}
\label{app:vla_pitfall}
Two Qwen3-VL-4B base VLMs we tried produced qualitatively different Path B behaviour. Our first attempt used an ``Instruct-Action'' variant that adds $2{,}048$ \texttt{<robot\_action\_*>} embeddings (initialised at the mean of the existing vocab, kept unfrozen). The OFT head we train does \emph{not} predict those tokens, but the rows for them sit in the tied \texttt{embed\_tokens} / \texttt{lm\_head} matrix and receive gradient from the retrofit's distillation loss and from label smoothing. Path B under this dirty base reaches $94.55\%$ 4-suite avg (Long-$10$ $84.8\%$); Path B under the clean base---no extra tokens added---reaches $96.70\%$ (Long-$10$ $94.2\%$). A $+9.4$pp Long-$10$ swing is entirely explained by whether $2{,}048$ unused rows in the shared matrix receive stray gradient. Use the clean base for retrofit-$+$-OFT; only add bespoke action vocabulary when the head actually predicts those tokens.

\subsection{VLA landscape: open VLAs on standard benchmarks}
\label{app:vla_landscape}
\begin{table}[h]
\caption{Open VLA landscape on standard benchmarks. LIBERO entries are success rates ($\%$) per task suite. Following common practice, $\pi_0$ and Octo-Base LIBERO numbers are taken from OpenVLA-OFT~\citep{kim2025openvlaoft}, which re-ran all three policies under a uniform LIBERO protocol. Our rows use the $30$k Path B numbers from Table~\ref{tab:vla_libero}.}
\label{tab:vla_landscape}
\centering\small
\setlength{\tabcolsep}{4pt}
\begin{tabular}{lccccccc}
\toprule
Model & \#Params & Spatial & Object & Goal & Long-10 & ref. \\
\midrule
Octo-Base    & 93M & 78.9 & 85.7 & 84.6 & 51.1 & \citep{team2024octo, kim2025openvlaoft} \\
OpenVLA      & 7B  & 84.7 & 88.4 & 79.2 & 53.7 & \citep{kim2024openvla, kim2025openvlaoft} \\
TraceVLA     & 7B  & 84.6 & 85.2 & 75.1 & 54.1 & \citep{zheng2024tracevla} \\
SpatialVLA   & 4B  & 88.2 & 89.9 & 78.6 & 55.5 & \citep{qu2025spatialvla} \\
$\pi_0$      & 3B  & \textbf{96.8} & \textbf{98.8} & 95.8 & 85.2 & \citep{black2024pi0, kim2025openvlaoft} \\
\midrule
\rowcolor{black!6}
\textbf{Qwen3-VL-2B + AR-Retrofit + OFT} (\emph{ours}) & 2B & $\mathbf{97.8}$ & $99.6$ & $\mathbf{98.6}$ & $92.6$ & this work \\
\rowcolor{black!6}
\textbf{Qwen3-VL-4B + AR-Retrofit + OFT} (\emph{ours}) & 4B & $94.6$ & $\mathbf{99.8}$ & $\mathbf{98.2}$ & $\mathbf{94.2}$ & this work \\
\bottomrule
\end{tabular}
\end{table}

\subsection{VLA \reskip{} setup: per-block drift and eligible-set selection}
\label{app:vla_pareto}

The \reskip{} eligible set $\mathcal{P}$ for each scale is selected from a per-block ablation \emph{on the trained Path B policy}, not on the VLM retrofit alone, because the VLA action stream activates a different routing distribution than the language stream. We measure per-block action-drift MSE on a $24$-sample sim trajectory: the policy is run with no skip and with each block individually skipped; the per-step action MSE between the skipped and no-skip rollouts is the per-block ``cost of removing this block on the action distribution''. The two safest blocks per scale form $\mathcal{P}$.

\textbf{2B Path B v2 30k} per-block action-drift MSE (sorted ascending, in units $\times 10^{-3}$):
block~$1$: $0.4$; block~$4$: $34.0$; block~$0$: $58.1$; block~$5$: $87.0$; remaining blocks $>$$100$. Selecting the two lowest: $\mathcal{P}_{2\text{B}}=\{1, 4\}$.

\textbf{4B Path B 30k clean} per-block action-drift MSE: block~$1$: $0.6$; block~$2$: $11.0$; block~$0$: $34.0$; block~$3$: $63.0$; remaining $>$$120$. Selecting: $\mathcal{P}_{4\text{B}}=\{1, 2\}$.

\paragraph{Cross-scale transfer of $\mathcal{P}$ fails.} Naively transferring $2$B's $\mathcal{P}{=}\{1, 4\}$ onto $4$B catastrophically drops LIBERO-Spatial to $12.4\%$ at $q{=}0.85$ — block-$4$'s drift on $4$B is roughly $10\times$ that of block-$2$, so the same eligible-set rule produces wildly different operating points across scales. $\mathcal{P}$ must be calibrated on each model's own sim distribution; this is the per-scale half of the modality-aware skip protocol.

\paragraph{Threshold calibration.} Per-block thresholds $\tau_n$ are the $q$-th empirical quantile of $w_{\text{recent},n}$ over $31{,}286$ ($2$B) / $31{,}257$ ($4$B) sim-rollout records (one record per decoded action token across many trajectories). The Tab.~\ref{tab:vla_reskip_pareto} numbers use this protocol. A separate ``Method A'' that calibrates on retrofit pretokenized data (as if the VLA were just a long-context LM) accidentally lands at a conservative operating point that mimics $q{=}0.99$ on the action distribution; we use it as a sensitivity comparator in our internal ablations and recommend the action-distribution-calibrated thresholds for any deployment.

\paragraph{Cross-modality threshold transfer is not safe.} The LAMBADA-calibrated thresholds at $q{=}0.85$ from Tab.~\ref{tab:retrofit_pareto}, applied unchanged to the same Path B policy at LIBERO inference, drop LIBERO-Spatial from $99.5\%$ (a $L{=}4$ $v3$ warm-start, partial eval) to $64\%$. The action distribution has more low-$w_{\text{recent}}$ mass than the language distribution at the same blocks; the language $q{=}0.85$ corresponds to roughly the action $q{=}0.50$, which over-skips. This is the empirical hard motivation for per-modality, per-token-class calibration in deployments that mix modalities within a single rollout.

\subsection{VLA planned follow-ups}
\label{app:vla_future}
\emph{Per-modality, per-token-class \reskip{}.} The Tab.~\ref{tab:vla_reskip_pareto} thresholds are calibrated on the action distribution as a whole; an obvious refinement is per-modality $(\mathcal{P}^{(m)}, \tau_n^{(m)}, M_{\max}^{(m)})$ within a single rollout (vision tokens, language tokens, action tokens). Working hypothesis: vision $\alpha$ concentrates on early blocks (allows aggressive late-block skip on perception steps); action $\alpha$ spreads to late blocks (restricts late-block skip on action prediction steps). The calibration harness (\texttt{retrofit/eval/calibrate\_vla\_thresholds.py}) is in place; the failed cross-modality transfer of Appendix~\ref{app:vla_pareto} is the empirical motivation. \emph{Compile $+$ skip composition.} A static-graph variant of the skip rule (compile-safe top-$k$ with a fixed routing decision per shape bucket) would let the iso-cost compile mode and the eager skip mode coexist in a single forward and stack their $1.029\times$ and $5$--$9\%$ savings. \emph{Edge hardware.} Wall-clock is H100 / bf16; RTX 4090 / Jetson Orin direct measurement is planned.

\section{NeurIPS Paper Checklist}
\label{app:checklist}
\begin{enumerate}
\item \textbf{Claims.} Yes. The abstract, introduction, and conclusions state the main claims and distinguish deployment-style max/min VLA reporting from mean-of-seeds sensitivity where repeated rollouts exist.
\item \textbf{Limitations.} Yes. Appendix~\ref{app:limitations} lists architecture coverage, compile / dynamic-skip composition, hardware scope, and per-modality calibration limits.
\item \textbf{Theory assumptions.} Not applicable. The paper is empirical and algorithmic; all equations define implemented objectives or routing rules.
\item \textbf{Experimental reproducibility.} Yes. Sections~\ref{sec:reskip}--\ref{sec:vla_results} and the appendix specify model scales, block partitions, training steps, datasets, calibration rules, and evaluation protocols.
\item \textbf{Open access to code and data.} Partially. Experiments use public datasets and benchmarks where available; release details for retrofit / VLA training scripts will follow the anonymous-submission policy.
\item \textbf{Compute.} Yes. Main text and appendices report GPU type, training steps, rough GPU-minutes / GPU-hours, and latency measurement settings.
\item \textbf{Dataset and benchmark provenance.} Yes. All public datasets and benchmarks are cited; LIBERO and lmms-eval protocols are described with rollout counts or split usage.
\item \textbf{Human subjects.} Not applicable. No new human-subject data are collected.
\item \textbf{Privacy.} Not applicable beyond the privacy considerations of the cited public datasets.
\item \textbf{Licenses.} Partially. Public resources are cited; final camera-ready release will include the exact code / model / data license table.
\item \textbf{Broader impacts.} Yes. The work targets lower inference cost and adaptive computation in pretrained models; no safety-critical deployment is claimed.
\item \textbf{Safeguards.} Yes. VLA thresholds are calibrated per model and action distribution; cross-modality transfer failures are explicitly documented as a deployment risk.
\end{enumerate}

\end{document}
